% Generated draft: requires mathematical and source review.
\documentclass[11pt]{article}
\usepackage[margin=25mm]{geometry}
\usepackage{fontspec}
\setmainfont{Latin Modern Roman}
\usepackage{amsmath,amssymb,mathtools,unicode-math}
\setmathfont{Latin Modern Math}
\usepackage{enumitem,needspace,xurl,hyperref}
\setlength{\parindent}{0pt}
\setlength{\parskip}{5pt}
\setlength{\emergencystretch}{3em}
\newcommand{\N}{\mathbb{N}}
\newcommand{\Z}{\mathbb{Z}}
\newcommand{\Q}{\mathbb{Q}}
\newcommand{\R}{\mathbb{R}}
\newcommand{\C}{\mathbb{C}}
\newcommand{\F}{\mathbb{F}}
\newcommand{\A}{\mathbb{A}}
\newcommand{\PP}{\mathbb{P}}
\newcommand{\E}{\mathbb{E}}
\newcommand{\eps}{\varepsilon}
\newcommand{\dd}{\,\mathrm{d}}
\begin{document}
\section*{UnsolvedMath ID: 800424}
\textbf{Source problem code:} catalogue entry
\subsection{Definitions and mathematical statement}
Let \(K\) be a number field, let \(f:\mathbb{P}^N_K\to\mathbb{P}^N_K\) be an endomorphism defined over \(K\) of degree at least \(2\), and let \(X\subsetneq\mathbb{P}^N_K\) be a subvariety defined over \(K\) and irreducible over \(K\). Geometric irreducibility of \(X\) is not assumed. Fix an algebraic closure \(\overline K\). Write \(f^n\) for the \(n\)-fold iterate of \(f\), and define the forward orbit of \(X\) by
\[
\mathcal O_f^+(X)=\bigcup_{n\geq 1}f^n(X).
\]
Its Zariski density means that \(\overline{\mathcal O_f^+(X)}^{\mathrm{Zar}}=\mathbb P^N_K\).

The supplied record does not define eventual stability or geometric eventual stability. The following is a provisional formulation using bounded numbers of irreducible components; agreement with the source's intended definitions requires review. For \(L\in\{K,\overline K\}\), let \(f_L\) and \(X_L\) denote base change to \(L\), and put
\[
r_L(n)=\#\left\{\text{irreducible components over }L\text{ of }\bigl((f_L^n)^{-1}(X_L)\bigr)_{\mathrm{red}}\right\},\qquad n\geq 1.
\]
Here the reduced inverse image is used, so components are counted without multiplicity. Provisionally, eventual stability over \(K\) means
\[
\exists B_K\in\mathbb N\ \forall n\geq 1,\qquad r_K(n)\leq B_K,
\]
and geometric eventual stability means
\[
\exists B_{\overline K}\in\mathbb N\ \forall n\geq 1,\qquad r_{\overline K}(n)\leq B_{\overline K}.
\]
Under this explicitly provisional interpretation, the supplied questions are:
\begin{enumerate}
\item[(a)] For every \(K,f,X\) as above, if \(\mathcal O_f^+(X)\) is Zariski dense in \(\mathbb P^N_K\), must \((f,X)\) be eventually stable over \(K\)?
\item[(b)] For every \(K,f,X\) as above, if \(\mathcal O_f^+(X)\) is Zariski dense and \(\dim X\geq 1\), must \((f,X)\) be geometrically eventually stable?
\item[(c)] For every \(K,f,X\) as above with \(\dim X\geq 1\), is eventual stability over \(K\) equivalent to geometric eventual stability? No Zariski-density hypothesis is imposed in this part.
\end{enumerate}
\subsection{Short English statement}
For an endomorphism of projective space of degree at least two and a proper subvariety irreducible over the number field of definition, ask whether a dense forward orbit forces eventual stability over that field, whether it forces geometric eventual stability in positive dimension, and whether the two stability notions are equivalent in positive dimension without assuming density. The component-count definitions used here are provisional because the supplied extract does not define these notions.
\subsection{Sources}
\begin{itemize}
\item Supplied local UnsolvedMath registry record, permanent ID \(800424\), exact source problem code \texttt{catalogue entry}, in the AMR Open Problem Lists collection. Registry title: \emph{Open Problems in Arithmetic Dynamics --- Question 47}.
\item Attributed in the supplied record to Benedetto, DeMarco, Ingram, Jones et al., \emph{Current Trends and Open Problems in Arithmetic Dynamics} (2018), Question 47. Supplied source URL: \url{https://arxiv.org/abs/1806.04980}.
\item Supplied provenance states \texttt{Accessed: 2026-10-16} and extraction method \texttt{arxiv\_latex}. This access information is reproduced from the record; no independent source visit or verification is asserted here.
\item The registry labels the problem \texttt{open}. This is a supplied status label, not an independently verified assessment of the literature.
\end{itemize}
\end{document}
